\SourceLeaf{386}%
the measurements themselves allow, so that from the results (the geodetic longitudes and latitudes) one can calculate back to the angle measurements with at least the accuracy with which they were made, the abbreviated formulas fall far short of what is needed. For very large triangles one would then even wish to be able to take the terms of the fourth order into account. With a different form of calculation this can be achieved by very flexible methods. This is no place to develop the matter further, and we content ourselves here with pointing out this need of higher geodesy. The work itself offers several cases in which the great advantages of such a treatment can be felt. For example, the intersection observations of the triangulation points Masse, Monte Soglio and Rocca Melone, made at the Turin observatory during the determinations of azimuth, are not used at all. On the assumption just stated, they could easily have provided a very valuable check and improved the accuracy.

